\section{总体架构与运行流程}
\subsection{系统边界与分层}
系统采用单仓组织智能体核心、QQ 插件、Web 控制台、MCP 服务和部署工具。Python 核心通过 Domain 数据类型与 Ports 接口表达消息、策略、模型、能力和输出，AstrBot 组合根负责装配具体实现。外部模型平台、校园认证系统和校园数据源提供服务；公共域名、统一登录与站点网关属于部署环境，不是智能体核心组件。

\begin{figure}[H]\centering
\begin{tikzpicture}[x=1cm,y=1cm]
\node[box,text width=6.1cm] (qq) at (0,0) {QQ 群成员\\NapCat / OneBot v11};
\node[box,text width=6.1cm] (web) at (7.1,0) {管理员 Web 工作台\\Vue 3 / Node.js Gateway};
\node[box,text width=6.1cm] (ad) at (0,-1.65) {AstrBot 输入与输出适配\\消息身份、引用、群历史、投递};
\node[box,text width=6.1cm] (pr) at (7.1,-1.65) {Runtime 预览与配置接口\\服务端同群历史、运行记录};
\node[box,text width=13.2cm,minimum height=2.1cm] (rt) at (3.55,-3.9) {\textbf{Dududa Agent Runtime}\\[2mm]上下文 $\rightarrow$ 感知 $\rightarrow$ 社交策略 $\rightarrow$ 模型路由\\[1mm]能力计划与执行 $\rightarrow$ 回答组织 $\rightarrow$ Persona $\rightarrow$ 输出请求};
\node[box,text width=3.75cm,minimum height=1.35cm] (llm) at (-1.05,-6.4) {模型适配器\\DeepSeek\\Chat Completions};
\node[box,text width=3.75cm,minimum height=1.35cm] (mcp) at (3.55,-6.4) {Unified MCP / 本地能力\\五类 MCP 校园服务\\校车时刻表};
\node[box,text width=3.75cm,minimum height=1.35cm] (state) at (8.15,-6.4) {策略与状态\\Persona / 群配置\\SQLite / JSON};
\node[future,text width=13.2cm] (future) at (3.55,-8.15) {独立模块与后续接入：Memory v2、主动日报、离线 Bandit、多步 Skill};
\draw[flow] (qq)--(ad);\draw[flow] (web)--(pr);
\draw[flow] (ad.south)--(ad.south |- rt.north);
\draw[flow] (pr.south)--(pr.south |- rt.north);
\draw[flow] (rt.south -| llm.north)--(llm.north);
\draw[flow] (rt)--(mcp);
\draw[flow] (rt.south -| state.north)--(state.north);
\draw[flow,dashed] (future.east) -- ++(0.4,0) |- (rt.east);
\end{tikzpicture}

\caption{嘟嘟哒的系统组成。实线为当前交互或调用关系；虚线模块为独立实现或后续接入方向。}
\label{fig:architecture}
\end{figure}

\begin{longtable}{@{}L{24mm}L{45mm}L{80mm}@{}}
\caption{架构分层与责任}\label{tab:layers}\\\toprule
层次 & 主要组件 & 责任与输入输出\\\midrule\endfirsthead
\toprule 层次 & 主要组件 & 责任与输入输出\\\midrule\endhead
交互与适配 & QQ、NapCat、AstrBot；Vue 3、Node.js & 将平台消息变为统一消息，将结果变为 QQ 消息或 Web 预览；控制台读取配置和运行记录。\\
领域与端口 & Domain、Ports & 表达 Actor、ConversationScope、MessageEnvelope、调用预算和版本化结果；不携带框架事件或密钥。\\
智能体编排 & Runtime、Perception、Social、Router & 管理阶段流转，构建上下文，识别意图，决定静默、澄清、回答或使用工具，选择计算配置。\\
业务能力 & Registry、Planner、Executor、Validator & 筛选本群可用能力，绑定输入，执行有界查询，将结果整理为 Observation。\\
内容与表达 & DirectChat、Composer、Persona、Final Validator & 使用问题或可信工具资料生成正文；组织结构、来源与篇幅；完成最终格式校验。\\
基础设施 & 模型适配器、Unified MCP、SQLite/JSON、Compose & 提供外部访问、服务会话、配置和持久状态；隔离第三方依赖及运行数据。\\\bottomrule
\end{longtable}

\subsection{核心组件与依赖方向}
组合根位于 \code{apps/astrbot-plugins/astrbot_plugin_dududa_core/}，其中 \code{composition.py} 创建规则与模型混合感知、确定性社交策略、静态模型路由、单步能力运行时、Persona Registry 及输出请求构建器。

核心 Runtime 的类名仍为 \code{OfflineRuntimeOrchestrator}，沿用了早期离线实现名称。当前正式组合已经为它注入真实模型和 MCP 适配器，因此类名中的 Offline 不代表正式运行使用假模型。测试则可以替换这些接口而保留相同编排。

\begin{center}
\begin{tikzpicture}[node distance=7mm]
\node[box,text width=133mm] (a) {应用组合根与平台适配器\quad AstrBot / Web / Model / MCP};
\node[box,text width=133mm,below=of a] (b) {应用服务与 Runtime\quad 上下文、感知、社交、路由、能力、回答};
\node[box,text width=133mm,below=of b] (c) {Domain 与 Ports\quad 统一数据模型、接口、调用与结果契约};
\draw[flow] (a)--(b);\draw[flow] (b)--(c);
\end{tikzpicture}
\end{center}
核心依赖抽象端口，应用组合根提供具体实现。新增平台需要适配消息身份、引用和投递回执；新增模型需要满足统一请求与结果协议；新增业务则实现能力定义和 Provider。这种拆分支持独立验证，但也要求维护类型、组合根与接口实现之间的一致性。

\subsection{关键数据对象}
\begin{tabularx}{\textwidth}{@{}L{38mm}Y@{}}\toprule
对象 & 业务含义\\\midrule
MessageEnvelope & 本轮平台消息、正文、时间、提及、引用及附件描述；作为运行输入。\\
Actor / ConversationScope & 当前行为主体及平台、机器人、会话、群的范围；决定可读取和可回复的上下文。\\
PerceptionResult & 意图、实体、引用关系、能力类别及置信信息；提供后续决策所需的结构化语义。\\
SocialDecision / ResponsePlan & 参与动作与理由、回答目标、模型档位和表达预算；区分行为选择与文本生成。\\
ToolPlan / Observation & 待执行的业务能力与参数，以及校验后的工具事实、来源和结果状态。\\
Draft / FinalResponse & 回答结构与最终正文；Persona 和格式校验在这一层完成。\\
DeliveryRequest / Receipt & 期望投递的目标、分片和内容，以及平台确认的实际结果。\\\bottomrule
\end{tabularx}

\subsection{一次查询的完整路径}
以“评课社区里数学分析 B1 的评价怎样”为例，系统首先从事件取得可信身份和同群上下文。规则识别明确的评课来源，快速档模型提取课程实体和意图；社交策略允许回应，能力目录筛选本群可用的评课能力；参数绑定器形成一次业务查询。MCP 返回课程、教师、评论和链接后，回答模型按事实组织文字，Composer 和 Persona 处理表达，最终由对应输出适配器完成呈现。

\begin{figure}[H]\centering
\begin{tikzpicture}[x=1cm,y=1cm]
\foreach \x/\name/\id in {0/入口适配/a,3.5/Runtime/b,7/模型适配/c,10.5/能力与 MCP/d}{
  \node[box,text width=2.6cm] (\id) at (\x,0) {\name};
  \draw[dashed,rulegray] (\x,-0.6)--(\x,-7.25);
}
\draw[flow] (0,-1)--node[above,font=\footnotesize]{统一消息与同群上下文}(3.5,-1);
\draw[flow] (3.5,-2)--node[above,font=\footnotesize]{结构化感知请求}(7,-2);
\draw[flow] (7,-2.7)--node[above,font=\footnotesize]{意图、实体与置信信息}(3.5,-2.7);
\draw[flow] (3.5,-3.6)--node[above,font=\footnotesize]{资格筛选、参数绑定、单步执行}(10.5,-3.6);
\draw[flow] (10.5,-4.4)--node[above,font=\footnotesize]{Observation：事实、身份、来源、状态}(3.5,-4.4);
\draw[flow] (3.5,-5.25)--node[above,font=\footnotesize]{资料与回答约束}(7,-5.25);
\draw[flow] (7,-5.95)--node[above,font=\footnotesize]{生成正文与用量}(3.5,-5.95);
\draw[flow] (3.5,-6.8)--node[above,font=\footnotesize]{格式校验后输出}(0,-6.8);
\end{tikzpicture}

\caption{校园查询的请求时序。模型承担感知与回答，工具调用由程序显式执行。}
\label{fig:sequence}
\end{figure}

一次典型查询由感知、业务能力执行和回答生成构成。业务能力内部可能读取多个数据条目，但当前 Planner 的计划仍只包含一个业务步骤；这与通用智能体跨服务自主规划多步任务不同。多项独立需求超出本轮范围时，当前实现会澄清或引导拆分。

\subsection{阶段状态与完成条件}
\begin{figure}[H]\centering
\begin{tikzpicture}[node distance=7mm and 7mm]
\node[box,text width=3.8cm] (r) {接收与预处理\\received / preprocessed};
\node[box,text width=3.8cm,right=of r] (c) {上下文与感知\\context\_ready / perceived};
\node[box,text width=3.8cm,right=of c] (d) {参与决策\\decided};
\node[box,text width=3.8cm,below=of d] (t) {计划与执行\\tools\_planned / executed};
\node[box,text width=3.8cm,left=of t] (v) {结果与内容\\validated / composed};
\node[box,text width=3.8cm,left=of v] (p) {表达与输出准备\\rendered / ready\_to\_emit};
\node[box,text width=13cm,below=of v] (end) {输出回执 $\rightarrow$ 完成\quad delivery\_acknowledged $\rightarrow$ completed};
\draw[flow] (r)--(c);\draw[flow](c)--(d);\draw[flow](d)--(t);\draw[flow](t)--(v);\draw[flow](v)--(p);
\draw[flow](p.south)--(p.south |- end.north);
\end{tikzpicture}

\caption{主要运行阶段。直接回答跳过工具阶段；错误和延期可从相应阶段终止。}
\end{figure}

运行状态记录从收到消息到生成最终内容的进展，输出请求与实际发送分开确认。\code{ready_to_emit} 只表示内容已准备好，\code{delivery_acknowledged} 表示收到输出适配器回执。Web 预览把结果交给预览适配器，不能据此推断 QQ 已投递。状态类型中保留 Memory 阶段，但当前生产组合未装配长期记忆读写。
\source{[S2, S3, S4]；源码：\nolinkurl{runtime/state.py}、\nolinkurl{runtime/orchestrator.py}、\nolinkurl{composition.py}。}
